\section{The Albert algebra: ten dimensions, a spinor, and matter parity}
\label{sec:lorentz}

Section~7 produced, from the clock, a $27$-dimensional Euclidean Jordan algebra. By the
Jordan--von Neumann--Wigner classification there is exactly one such algebra that is
``exceptional'': the \emph{Albert algebra} $H_3(\mathbb O)$ of $3\times3$ Hermitian matrices
over the octonions. Since the 1930s it has been suspected of having something to do with
quantum mechanics and particle physics; this section computes its symmetries on the
committed basis and finds the Lorentz group of ten-dimensional space-time, a chiral spinor,
and the repository's own matter parity.

\begin{plainwords}
The \emph{octonions} are the largest of the four number systems in which you can divide:
the real numbers, the complex numbers, the quaternions, and the octonions, of dimensions
$1,2,4,8$. Octonions do not even obey $(ab)c=a(bc)$, yet $3\times3$ ``Hermitian'' matrices of
them --- real numbers on the diagonal, octonions above it, their conjugates below --- form a
consistent algebra of dimension $3+3\times8=27$. Its symmetries are the exceptional groups
$F_4$ and $E_6$ \cite{Baez2002}. Manogue, Dray and Wilson have recently described the
Standard Model through octonionic $E_8$ structures of this kind \cite{MDW2022}.
\end{plainwords}

\subsection{Anatomy of a \texorpdfstring{$3\times3$}{3x3} octonionic matrix}

Choose the idempotent $c=E_{11}$ (a $1$ in the top-left corner). Every element of the
algebra splits by how it behaves under multiplication by $c$ --- its \emph{Peirce
decomposition}, with eigenvalues $1$, $\tfrac12$, $0$:

\begin{figure}[htbp]
\centering
\begin{tikzpicture}[scale=1.15,
   cellbox/.style={minimum width=15mm,minimum height=11mm,draw=white,line width=1.5pt,font=\small,align=center}]
  \node[cellbox,fill=ink,text=white] at (0,0) {$a$\\[-2pt]\tiny real};
  \node[cellbox,fill=gold!45] at (1.35,0) {$z$\\[-2pt]\tiny $\mathbb O$};
  \node[cellbox,fill=gold!45] at (2.7,0) {$\bar y$\\[-2pt]\tiny $\mathbb O$};
  \node[cellbox,fill=gold!45] at (0,-1.0) {$\bar z$};
  \node[cellbox,fill=teal!40] at (1.35,-1.0) {$b$\\[-2pt]\tiny real};
  \node[cellbox,fill=teal!40] at (2.7,-1.0) {$x$\\[-2pt]\tiny $\mathbb O$};
  \node[cellbox,fill=gold!45] at (0,-2.0) {$y$};
  \node[cellbox,fill=teal!40] at (1.35,-2.0) {$\bar x$};
  \node[cellbox,fill=teal!40] at (2.7,-2.0) {$c'$\\[-2pt]\tiny real};
  \draw[ink!60,line width=1pt] (-0.85,0.62) -- (-0.95,0.62) -- (-0.95,-2.62) -- (-0.85,-2.62);
  \draw[ink!60,line width=1pt] (3.55,0.62) -- (3.65,0.62) -- (3.65,-2.62) -- (3.55,-2.62);
  % legend
  \begin{scope}[xshift=5.1cm,yshift=0.1cm]
    \fill[ink] (0,0) rectangle (0.4,0.3); \node[anchor=west,font=\small] at (0.55,0.15) {$A_1$: \ $1$ dim --- the line through $c$};
    \fill[gold!45] (0,-0.8) rectangle (0.4,-0.5); \node[anchor=west,font=\small] at (0.55,-0.65) {$A_{1/2}$: \ $16=\mathbb O^2$ --- a spinor};
    \fill[teal!40] (0,-1.6) rectangle (0.4,-1.3); \node[anchor=west,font=\small] at (0.55,-1.45) {$A_0$: \ $10=H_2(\mathbb O)$ --- a vector};
    \node[anchor=west,font=\scriptsize,text=slate,align=left] at (0,-2.35)
      {$\det\begin{psmallmatrix} b & x\\ \bar x & c'\end{psmallmatrix}=bc'-|x|^2$\\
       has signature $(1,9)$: ten-dimensional\\Minkowski space};
  \end{scope}
\end{tikzpicture}
\caption{The Peirce decomposition $27=1+16+10$ of the Albert algebra with respect to a
primitive idempotent. The lower-right $2\times2$ block is a copy of ten-dimensional
Minkowski space; the two octonions in the first row and column form a sixteen-component
spinor.}
\label{fig:peirce}
\end{figure}

That the $2\times2$ Hermitian octonionic matrices form ten-dimensional Minkowski space, with
the determinant as the Minkowski interval, is a classical observation (Sudbery; Kugo and
Townsend; Baez and Huerta), and it is the algebraic reason the superstring lives in ten
dimensions \cite{BaezHuerta,KugoTownsend}. The new content is that the forty-point clock
reaches it, on the committed basis.

\subsection{The symmetries, computed}

\begin{result}{Theorem 8.1 (the clock Albert algebra) \hfill \tierC}
For the $27$-dimensional clock algebra $A$ of Theorem~7.2:
\begin{enumerate}[label=(\roman*),leftmargin=1.5em]
\item the derivations $\mathrm{Der}(A)=\mathrm{span}\,[L_x,L_y]$ form a $52$-dimensional
      Lie algebra with negative-definite trace form $(0,52)$: the compact $F_{4(-52)}$;
\item the reduced structure algebra $\mathrm{str}_0(A)=\mathrm{Der}(A)\oplus L(A_{\mathrm{tr}=0})$
      is $78$-dimensional, obeys $[D,L_x]=L_{Dx}$ for all $52\times27$ pairs, and has
      inertia $(26,52)$: $E_{6(-26)}$;
\item the derivations fixing $c$ form a compact $36$-dimensional $\mathfrak{spin}(9)$;
      adding the nine ``boosts'' $L_y$ ($y\in A_0$, orthogonal to $e-c$) gives a
      $45$-dimensional algebra that preserves $A_0$ and its determinant of inertia $(1,9)$,
      and acts faithfully: it is $\so(1,9)$;
\item on $A_{1/2}$ the commutant of $\so(1,9)$ is one-dimensional: the $16$ is absolutely
      irreducible --- the Majorana--Weyl spinor of $\so(1,9)$.
\end{enumerate}
\cert{w33\_pass10950\_clock\_albert\_lorentz\_spinor.py}
\end{result}

\begin{center}
\begin{tikzpicture}[x=3.25cm,y=0.95cm,
   hd/.style={font=\sffamily\scriptsize\bfseries,align=center,text=ink},
   rw/.style={font=\sffamily\scriptsize,align=right,anchor=east,text=slate},
   cl/.style={draw=#1,fill=#1!10,rounded corners=3pt,minimum width=29mm,minimum height=7.2mm,font=\small,align=center},
   ok/.style={font=\scriptsize\bfseries,text=teal!80!black}]
  \node[hd] at (1,0.95) {ordinary space-time\\$H_2(\C)$: \ $4$ dims};
  \node[hd] at (2,0.95) {superstring space-time\\$H_2(\mathbb O)$: \ $10$ dims};
  \node[hd] at (3,0.95) {exceptional space-time\\$H_3(\mathbb O)$: \ $27$ dims};
  \foreach \lab/\y in {{rotations (automorphisms)}/0,{Lorentz (norm-preserving)}/-1,{conformal}/-2,{quasiconformal}/-3}{
     \node[rw] at (0.45,\y) {\lab};}
  % column 1
  \node[cl=slate] at (1,0) {$SO(3)$};
  \node[cl=slate] at (1,-1) {$SO(1,3)$};
  \node[cl=slate] at (1,-2) {$SO(2,4)$};
  \node[font=\scriptsize,text=slate!60] at (1,-3) {---};
  % column 2
  \node[cl=sage] at (2,0) {$\mathrm{Spin}(9)$ \ \computed};
  \node[cl=sage] at (2,-1) {$\mathrm{Spin}(1,9)$ \ \computed};
  \node[cl=slate] at (2,-2) {$SO(2,10)$};
  \node[font=\scriptsize,text=slate!60] at (2,-3) {---};
  % column 3
  \node[cl=teal] at (3,0) {$F_{4(-52)}$ \ \computed};
  \node[cl=teal] at (3,-1) {$E_{6(-26)}$ \ \computed};
  \node[cl=teal] at (3,-2) {$E_{7(-25)}$ \ \computed};
  \node[cl=ink,text=white,fill=ink] at (3,-3) {$E_{8(-24)}$ \ \computed};
  \node[font=\sffamily\scriptsize,text=slate,align=left,anchor=west] at (0.55,-4.0)
    {\computed\ \ computed in this paper on the committed $E_8$ basis, in the real form selected by the Hesse clock\\
     \phantom{\computed}\ \ (Passes 10949--10950); unmarked cells are classical.};
\end{tikzpicture}
\captionof{figure}{The ladder of space-times. A Jordan algebra of Hermitian matrices is a
space-time; its determinant is the interval. For $2\times2$ complex matrices it is
ordinary Minkowski space, for $2\times2$ octonionic matrices the ten-dimensional space of the
superstring, and for $3\times3$ octonionic matrices the $27$-dimensional ``exceptional
space-time'' of G\"unaydin \cite{Gunaydin1993}. Each has rotations, a Lorentz group, a conformal
group and (for the last) a quasiconformal group. The Hesse clock's real form of $E_8$ produces
the whole right-hand column, and the middle column inside it.}
\label{fig:spacetimes}
\end{center}

\begin{physical}[What it means physically: the clock picks the division algebra]
The last row of Freudenthal's ``magic square'' of Lie algebras can be built from the octonions
in two ways: from the division algebra $\mathbb O$, giving the chain
$F_{4(-52)}\subset E_{6(-26)}\subset E_{7(-25)}\subset E_{8(-24)}$ of automorphism, Lorentz,
conformal and quasiconformal groups of $H_3(\mathbb O)$, or from the split octonions, giving
the split forms. The Hesse clock selects the first: the \emph{division-algebra} octonions, the
ones with a positive norm, the ones that make $H_2(\mathbb O)$ a genuine Minkowski space. The
repository's earlier ``split'' choice lands on the other side. \tierC\ (the four real forms)
\ \tierB\ (the reading)
\end{physical}

\subsection{Down to four dimensions}

Choose a further splitting of the $10$ into $4+6$: time $u=e-c$, one frame direction, and
a complex line inside the octonions supplied by the complex structure of $\mathfrak p^+$
itself. The subalgebra of $\so(1,9)$ preserving the $4$ and the $6$ is $21$-dimensional,
$\so(1,3)\oplus\so(6)\cong\mathfrak{sl}(2,\C)\oplus\mathfrak{su}(4)$, and its commutant on
the $16$ is two-dimensional and contains a complex structure. Hence
\[
\mathbf{16}\;=\;(\mathbf2,\mathbf4)\;\oplus\;\overline{(\mathbf2,\mathbf4)} .
\]
\tierC\ (\cert{w33\_pass10950\_clock\_albert\_lorentz\_spinor.py})

\begin{physical}[Two readings of the same branching]
A ten-dimensional chiral spinor, seen by a four-dimensional observer, becomes four
left-handed Weyl spinors transforming as the $\mathbf4$ of $SU(4)$. If $SU(4)$ is read as
the rotation group of six hidden dimensions, this is the fermion content of maximally
supersymmetric Yang--Mills theory. If it is read as Pati--Salam colour, with lepton number
as the fourth colour, it is one generation of quarks and leptons. The branching is exact;
which reading nature uses is not decided by it. \tierB
\end{physical}

\subsection{Matter parity is a rotation by \texorpdfstring{$2\pi$}{2 pi}}

The repository carries an independent certificate for a ``matter charge'' $Q_\psi$ on the
$27$: it splits $27=16_1+10_{-2}+1_4$, and on the $45$ triads the charges fall into the
pattern $40\times(-2,1,1)+5\times(-2,-2,4)$. Matter parity is $(-1)^{Q_\psi}$: odd on the
$16$, even on the $10$ and $1$. In the Minimal Supersymmetric Standard Model a parity of
exactly this kind is what forbids the proton from decaying through dimension-four
operators.

\begin{result}{Theorem 8.2 (matter parity is the Peirce symmetry) \hfill \tierC}
For every one of the $27$ coordinate idempotents $c$, the Peirce charge
\[
Q=6L_c-2\qquad(\text{eigenvalue }4\text{ on }A_1,\;1\text{ on }A_{1/2},\;-2\text{ on }A_0)
\]
reproduces the certified spectrum of $Q_\psi$, is conserved on all $45$ triads, and
gives exactly the pattern $40\times(-2,1,1)+5\times(-2,-2,4)$. Its parity $(-1)^Q$ is the
Peirce symmetry $U_{2c-e}$, an automorphism of $A$ that is $+1$ on $A_1\oplus A_0$ and $-1$
on the spinor $A_{1/2}$. For a single-plane rotation generator $D\in\mathfrak{spin}(9)$,
$\exp(2\pi D)=U_{2c-e}$ (numerical agreement $2.7\times10^{-15}$).
\cert{w33\_pass10950\_clock\_albert\_lorentz\_spinor.py}
\end{result}

\begin{figure}[htbp]
\centering
\begin{tikzpicture}[>={Stealth[length=2.2mm]}]
  % charges
  \begin{scope}
    \foreach \lab/\q/\n/\c/\s [count=\i from 0] in {{$A_1$}/4/1/ink/{+},{$A_{1/2}$}/1/16/gold/{-},{$A_0$}/-2/10/teal/{+}}{
      \node[draw=\c,fill=\c!15,rounded corners=3pt,minimum width=22mm,minimum height=9mm,font=\small,align=center] at (0,-\i*1.15) {\lab: $\n$ dims};
      \node[font=\small] at (2.1,-\i*1.15) {$Q=\q$};
      \node[circle,draw=\c,fill=white,font=\small\bfseries,minimum size=6mm,inner sep=0pt] at (3.3,-\i*1.15) {$\s$};}
    \node[font=\sffamily\scriptsize,text=slate] at (3.3,0.75) {$(-1)^Q$};
  \end{scope}
  % 2pi rotation
  \begin{scope}[xshift=7.4cm,yshift=-1.15cm]
    \draw[->,rose,line width=1.3pt] (20:1.05) arc[start angle=20,end angle=375,radius=1.05];
    \node[font=\small,text=rose!80!black] at (0,1.4) {rotate by $2\pi$};
    \draw[->,teal,line width=1.2pt] (0,0) -- (0.8,0.35);
    \node[font=\scriptsize,text=teal!80!black,align=center] at (0,-1.55) {vectors and scalars: $+1$};
    \node[font=\scriptsize,text=gold!70!black,align=center] at (0,-1.95) {spinors: $-1$ \ (the $16$)};
  \end{scope}
  \node[font=\sffamily\small] at (4.95,-1.15) {$=$};
\end{tikzpicture}
\caption{Matter parity in the clock Albert algebra: the sign that a full turn gives a spinor.}
\label{fig:parity}
\end{figure}

\begin{physical}[What it means physically]
A rotation by $2\pi$ returns every vector to itself, but multiplies every spinor by $-1$:
the mathematical root of the difference between fermions and bosons. In $SO(10)$ grand
unification it is known that matter parity can be realised by exactly such a central
element of $\mathrm{Spin}(10)$, acting as $-1$ on the matter $\mathbf{16}$ and $+1$ on the
Higgs $\mathbf{10}$ (cf.\ \cite{Martin1992}). Theorem~8.2 is the objectwise version of that fact
inside the forty-point clock: the repository's matter parity \emph{is} the $2\pi$ rotation
of the $\mathrm{Spin}(9)$ fixing a primitive idempotent. The same $27=1+10+16$ that a
GUT reads as singlet, Higgs and matter is, in the clock real form, scalar, Lorentz vector
and Weyl spinor. \tierB\ (reading) \ \tierC\ (identity)
\end{physical}

\subsection{The clock's fourth tick is the same sign}

Section~4 met a clock whose fourth tick is a sign flip. The other research track has since
shown that this is not an analogy but one spin structure, seen at both ends of the ladder.
We cite their results (certificates named) rather than re-derive them.

\begin{result}{Theorem 8.3 (clock, spin and matter parity; Passes 10951--10959) \hfill \tierC}
\begin{enumerate}[label=(\roman*),leftmargin=1.5em]
\item The negacyclic clock tick lifts uniquely to $g=\begin{psmallmatrix}0&1\\1&1\end{psmallmatrix}\in GL(2,3)$
      of order eight, with $g^4=-I$; under an explicit isomorphism $SL(2,3)\cong 2T\subset
      \mathrm{Spin}(3)$ the element $-I$ is the quaternion $-1$, the central $2\pi$ spin sign.
      \cert{w33\_pass10951\_clock\_pin\_spin\_central\_sign\_bridge.py}
\item As a signed permutation of the four clock coordinates, determinant-odd ticks exchange
      the two $D_4$ half-spin sheets and determinant-even ticks preserve them.
      \cert{w33\_pass10955\_d4\_halfspin\_clock\_bridge.py}
\item On the Albert Peirce $16$, the $\mathrm{Spin}(8)$ fixing a spatial axis splits it as
      $8_s+8_c$; an exact rational $\mathrm{Spin}(9)$ half-turn $U$ exchanges them, with
      $U^2=-I_{16}$ --- the matter-parity sign of Theorem~8.2 --- and $U^4=I_{16}$.
      \cert{w33\_pass10956\_albert\_spin8\_halfspin\_normalizer.py},
      \cert{w33\_pass10958\_exact\_albert\_mu4\_missing\_mode\_intertwiner.py}
\item One Albert $16$ cannot carry the whole order-eight clock as a $GL(2,3)$ representation;
      the minimal carrier is $16\oplus16$.
      \cert{w33\_pass10959\_albert\_clock\_extension\_obstruction.py}
\end{enumerate}
\end{result}

\begin{center}
\begin{tikzpicture}[>={Stealth[length=2mm]},
   bx/.style={draw=#1,fill=#1!10,rounded corners=3pt,align=center,font=\sffamily\scriptsize,minimum height=10mm,minimum width=30mm}]
  \node[bx=teal] (g) at (0,0) {clock tick $g$\\order $8$};
  \node[bx=ink] (m) at (4.6,0) {$g^4=-I$\\central $2\pi$ sign};
  \node[bx=gold] (u) at (9.2,0) {Albert half-turn $U$\\$U^2=-I_{16}$};
  \node[bx=rose] (p) at (4.6,-1.6) {matter parity\\(Theorem 8.2)};
  \draw[->,thick] (g) -- node[above,font=\tiny]{four ticks} (m);
  \draw[->,thick] (u) -- node[above,font=\tiny]{two half-turns} (m);
  \draw[->,thick] (m) -- (p);
\end{tikzpicture}
\end{center}

\begin{boundary}
The $\so(1,9)$ here is the stabiliser of an idempotent inside $E_{6(-26)}$, an
\emph{internal} symmetry of the algebra; it is not derived to be the Lorentz group of
physical space-time, and the $2\pi$ rotation is a rotation in that internal $\mathrm{Spin}(9)$.
The identification of $Q_\psi$ with $6L_c-2$ holds up to the $W(E_6)$-transitive choice of
coordinate idempotent. Neither chirality selection, masses nor gauge dynamics follow.
\end{boundary}

\begin{keyidea}
The exceptional ladder closes: from one qutrit, through $W(3,3)$, $E_8$, a Hesse clock and
its real form, to the Albert algebra $H_3(\mathbb O)$ --- whose automorphisms are $F_4$,
whose structure group is $E_{6(-26)}$, and inside which sit ten-dimensional Minkowski space,
its Lorentz algebra $\so(1,9)$, a Majorana--Weyl spinor, and, as a rotation by $2\pi$, the
matter parity that protects the proton.
\end{keyidea}
